\documentclass[11pt]{article}
\usepackage{amsmath,amssymb,amsthm,hyperref}
\newtheorem{theorem}{Theorem}
\newtheorem{lemma}{Lemma}
\newcommand{\E}{\mathbb E}
\newcommand{\Pol}{\operatorname{Pol}}
\title{Quantitative relative balance at the endpoint}
\author{Asymptotic research note; independently reviewed}
\date{30 September 2026}
\setlength{\emergencystretch}{2em}
\begin{document}\maketitle
Assume $p_1,\ldots,p_m:[K]\to[0,1]$ are nondecreasing, the positive
row weights $w$ sum to one, and
\[
 \E_w\prod_{j\in S}p_j=\frac1{|S|+1}\qquad(S\subseteq[m]).
\]
Write $u_i=1-p_i$ and $s=\sum_i u_i$. Constants below are absolute,
independent of the number of rows, weights, and chosen column.

\begin{theorem}\label{main}
For all $m\ge2$ and every row,
\[
 \max_i(1-p_i)\le\frac{C(s+m^{-1/3})}{m}.           \tag{1}
\]
This strengthens the previously proved bound $C(s+1)/m$.
It is a structural tensor assertion; no new per-die exponent is
claimed in this note.
\end{theorem}

Fix a column $i$, put $M=m-1$, and write
$s_-=\sum_{j\ne i}u_j$. Set
\[
 \ell=\lfloor M/4\rfloor,\quad N=M-\ell,\quad
 \alpha=\ell/M,\quad x=\alpha s_-,\quad\theta=1/16.
\]
We work with sufficiently large $M$; smaller values are covered by
increasing $C$, since $u_i\le s$.
For a uniformly selected $\ell$-subset $B$ of the columns other than
$i$, put $A_B=\prod_{j\in B}p_j$ and $\overline A=\E_B A_B$.
There are two probability measures on the original rows:
\[
 \mu_0=(\ell+1)w\overline A,\qquad
 \mu_1=(\ell+1)(\ell+2)w u_i\overline A.
 \tag{2}
\]
Their exact density relation, including zero entries, is
\[
 d\mu_1=(\ell+2)u_i\,d\mu_0.                       \tag{3}
\]
The corresponding scalar probability densities are
\begin{align*}
 b_0(z)&=\frac{\ell+1}{\ell}(1-z/\ell)^\ell
                     \boldsymbol1_{[0,\ell]}(z),\\
 b_1(z)&=\frac{(\ell+1)(\ell+2)}{\ell^2}
                   z(1-z/\ell)^\ell\boldsymbol1_{[0,\ell]}(z).
\end{align*}
Let $\beta_a$ denote integration against $b_a$.

\begin{lemma}[The marked comparison]\label{marked}
For $a=0,1$, $1\le D\le\sqrt M$, and
$Q_b(z)=\sum_{j=0}^D c_{bj}L_j(\theta z)$ with
$B_b=\sum_j|c_{bj}|$, one has
\[
 \left|\E_{\mu_a}Q_1(x)Q_2(x)-\beta_a(Q_1Q_2)\right|
 \le C B_1B_2\left(\frac D M+e^{-cM}\right).        \tag{4}
\]
The constants are uniform in $i$.
\end{lemma}
\begin{proof}
For $a=0$, this is the scalar endpoint comparison in
\path{../../power_per_die_lower_bound.tex}, applied to the $M$
columns other than $i$. They still satisfy exactly the same uniform
mixed-moment identities with the original row weights. Thus its
application is not circular.

We detail the changes needed for $a=1$. The finite global balance
bound for all $m$ columns gives
\[
 u_i\le\frac{256(s_-+u_i+1)}{M+1}
       \le\frac{C(s_-+1)}M                         \tag{5}
\]
for large $M$, after moving the $u_i$ term to the left.
The finite balance and row-mass bounds applied to the remaining $M$
columns give
\[
 \max_{j\ne i}u_j\le C(s_-+1)/M,\quad
 \overline A\le e^{-\alpha s_-},\quad
 \E_w(s_-+1)^k e^{-c s_-}\le C_c^{k+1}k!/M.        \tag{6}
\]
In particular the marked normalization obeys, for integers $k\ge0$,
\[
 (\ell+1)(\ell+2)\E_w[u_i(s_-+1)^k e^{-c s_-}]
       \le C_c^{k+2}(k+1)!.                       \tag{7}
\]

For $P=Q_1Q_2$, multilinearity and the mixed moments give the exact
identity
\[
 (\ell+1)(\ell+2)\E_w\E_B
 [u_i A_B\Pol_N(P)((\ell u_j)_{j\notin B\cup\{i\}})]
       =\beta_1(P).                              \tag{8}
\]
Indeed its scalar integral is
$(\ell+1)(\ell+2)\int_0^1(1-t)t^\ell P(\ell(1-t))\,dt$.
The supports of the mark, forced successes, and polarized factors
are disjoint.

The pointwise polarization argument from the cited comparison proof
is unchanged. If $v$ is the mean of the remaining $\ell u_j$,
its centered variance is at most $CN(s_-+1)v$. The off-diagonal
Laguerre derivative estimate cancels the factor $v$ in the centered
coefficient estimate, giving pointwise error at most
\[
 B_1B_2e^{\theta v}
  \sum_{k=2}^{2D}\frac{(C\sqrt{Dk(s_-+1)/M})^k}{k!}.
\]
Average with $u_i A_B$, apply (7), and include its normalization.
The resulting sum is at most
\[
 C B_1B_2\sum_{k=2}^{2D}
 (C\sqrt{D/M})^k
 \frac{k^{k/2}(\lceil k/2\rceil+1)!}{k!}
 \le C B_1B_2 D/M.                               \tag{9}
\]
The extra factorial factor compared with the unmarked argument is
only polynomial in $k$; the geometric ratio tends to zero uniformly
for $D\le\sqrt M$.

To replace $v$ by $x$, use the same success-weighted subset law.
On $s_-\le M/256$ its bias and second moment satisfy
\begin{align*}
 0\le\E_B(v-x)&\le C s_-(s_-+1)/M,\\
 \E_B(v-x)^2&\le C\left(\frac{s_-(s_-+1)}M
                    +\frac{s_-^2(s_-+1)^2}{M^2}\right).
\end{align*}
These depend only on the remaining columns. The singular Taylor
remainder established in
\path{weighted_success_subset_moments.tex} bounds the pointwise
error by
\[
 C B_1B_2 e^{\theta\alpha s_-/(1-\alpha)}
 \left[\frac{\sqrt D\sqrt{s_-}(s_-+1)}M
       +\frac{D(s_-+1)}M+\frac{D s_-(s_-+1)^2}{M^2}\right].
\]
Multiplication by $u_i\overline A$ and (7) bound its normalized
integral by $CB_1B_2D/M$. On $s_->M/256$, the original uniform-subset
Cauchy--Schwarz bound leaves an exponentially decreasing factor;
the additional normalization $O(M^2)$ and $u_i\le1$ are absorbed
in $e^{-cM}$. This proves (4).
\end{proof}

\begin{lemma}[Laplace tests without a logarithmic loss]\label{laplace}
Uniformly for $1/2\le t\le M^{1/3}$ and $a=0,1$,
\[
 \left|\E_{\mu_a}e^{-t x}-\beta_a(e^{-t z})\right|
 \le C t/M+C\exp(-c\sqrt M/t).                   \tag{10}
\]
Consequently, for all sufficiently large $M$ and $1\le t\le M^{1/3}$,
\[
 c/t\le\E_{\mu_0}e^{-t x}\le C/t,\qquad
 \E_{\mu_1}e^{-t x}\le C/t^2.                    \tag{11}
\]
The upper bound for $\mu_0$ also holds when $t\ge1/2$.
\end{lemma}
\begin{proof}
The Laguerre generating function gives
\[
 e^{-tz/2}=(1-r)\sum_{j=0}^{\infty}r^jL_j(\theta z),
 \qquad r=\frac{t}{t+2\theta}.
\]
Truncate at $D=\lfloor\sqrt M\rfloor$, calling the resulting
polynomial $Q_D$. Since $|L_j(\theta z)|\le e^{\theta z/2}$,
\[
 |Q_D(z)-e^{-tz/2}|\le r^{D+1}e^{\theta z/2},\quad
 |Q_D(z)^2-e^{-tz}|\le2r^{D+1}e^{\theta z}.        \tag{12}
\]
Both $\mu_a$ and $\beta_a$ integrate $e^{\theta z}$ to at most an
absolute constant: this follows from (6)--(7) and
$\overline A\le e^{-\alpha s_-}$, or directly from their scalar
densities.

Apply (4) separately to $L_jL_k$, using degree $\max(j,k)$; the
constant product has zero error since both measures have mass one.
For $c_j=(1-r)r^j$, this yields
\begin{align*}
 |\E_{\mu_a}Q_D(x)^2-\beta_a(Q_D^2)|
 &\le\frac C M\sum_{j,k\le D}c_jc_k\max(j,k)+Ce^{-cM}\\
 &\le\frac{2C}M\sum_{j\ge0}j c_j+Ce^{-cM}\\
 &\le Ct/M+Ce^{-cM}.
\end{align*}
Furthermore $r^{D+1}\le e^{-c\sqrt M/t}$ for $t\ge1/2$.
Together with (12) this proves (10).
The scalar bounds are immediate from $b_0(z)\le2e^{-z}$,
$b_1(z)\le Cze^{-z}$, and $b_0(z)\ge1/2$ on $0\le z\le1/2$.
For $t\le M^{1/3}$, the error in (10) is negligible compared with
$1/t$, and at most $C/t^2$. This proves (11).
\end{proof}

\begin{lemma}[An interval with controlled density ratio]\label{interval}
There are absolute constants $0<a<b$ and $c,C>0$ such that, for
all large $M$ and $1\le t\le M^{1/3}$,
\[
 \mu_0(a/t<x<b/t)\ge c/t,\qquad
 \mu_1(x\le b/t)\le C/t^2.                       \tag{13}
\]
\end{lemma}
\begin{proof}
First the unmarked comparison (4) implies
\[
 \mu_0(x\le y)\le C(y+M^{-1/2})\qquad(y\ge0).    \tag{14}
\]
To see this directly, normalize
$T_D(\theta z)=\sum_{j<D}(1-j/D)L_j(\theta z)$ by its coefficient
sum $(D+1)/2$. Its square is bounded below by an absolute positive
constant on $z\le c_0/D$, has scalar expectation $O(1/D)$,
and has comparison error $O(D/M)$. Choose
$D$ comparable to $\min(\sqrt M,c_0/y)$; larger $y$ are covered by
probability at most one. This proves (14) using only the $M$-column
unmarked tensor.

By (11), $\E_{\mu_0}e^{-t x}\ge c_1/t$. Choose $a>0$ sufficiently
small that the bound (14) for $y=a/t$ is at most $c_1/(4t)$ for
large $M$. This is uniform because $tM^{-1/2}\le M^{-1/6}$.
For the upper tail,
\[
 \E_{\mu_0}[e^{-t x}\boldsymbol1_{x\ge b/t}]
 \le e^{-b/2}\E_{\mu_0}e^{-t x/2}\le Ce^{-b/2}/t.
\]
Choose $b$ large enough to make this at most $c_1/(4t)$.
The remaining interval therefore carries at least $c_1/(2t)$
weighted mass, hence at least that much probability. The marked
bound follows from (11) and
$\boldsymbol1_{x\le b/t}\le e^b e^{-t x}$.
\end{proof}

\begin{proof}[Proof of Theorem~\ref{main}]
Fix any row $q$. If $x(q)\le a/(2t)$, every row in the interval
of (13) has strictly larger $x$, and hence lies strictly before $q$.
Monotonicity gives $u_i\ge u_i(q)$ throughout that interval.
The density relation (3) and (13) give
\[
 (\ell+2)u_i(q)\frac c t
 \le\mu_1(a/t<x<b/t)\le C/t^2,
 \qquad u_i(q)\le C/(Mt).                        \tag{15}
\]
If $x(q)\le a/2$, take
$t=\min(M^{1/3},a/(2x(q)))$, with the second entry interpreted as
$+\infty$ at $x(q)=0$. This belongs to $[1,M^{1/3}]$, and (15) gives
\[
 M u_i(q)\le C(x(q)+M^{-1/3})\le C(s(q)+M^{-1/3}).
\]
If $x(q)>a/2$, then $s(q)\ge s_-(q)\ge a/(2\alpha)$ is bounded
below by a positive constant. The old global bound
$u_i\le256(s+1)/m$ is therefore already $u_i\le C s/m$.
The column $i$ was arbitrary, and $M=m-1$ is comparable to $m$.
This proves (1), including rows containing zero entries.
\end{proof}
\end{document}
